\section{Model}
\label{sec:method}

\subsection{Architecture}
\label{sec:method:arch}
The model maps a window of $T=20$ daily cross-sectionally normalised feature
vectors $\mathbf{x}_{1:T} \in \mathbb{R}^{T \times F}$, $F=21$, to a scalar
ranking score in three stages.

\paragraph{Temporal encoder.}
A dilated causal convolutional stack reduces the window to a vector
$\mathbf{e} \in \mathbb{R}^{n}$, where $n$ is the qubit count. Two blocks with
kernel size 3 and dilations $1$ and $2$ give a receptive field of seven trading
days; the representation of the final step is projected through a two-layer
head whose $\tanh$ output is scaled by $\pi$, so that every component is a valid
rotation angle. Padding is applied on the left only. A symmetric padding, as
used in the reference implementation this work audits, mixes each position with
the step that follows it and so does not implement the receptive field the
architecture claims.

\paragraph{Quantum convolutional layer.}
The encoder output is embedded as single-qubit rotations,
$\ket{\psi_0} = \bigotimes_{i=1}^{n} R_Y(e_i)\ket{0}$, and processed by
$L = \min(L_{\max}, \lfloor \log_2 n \rfloor)$ convolution-and-pooling stages.
Each stage applies the two-qubit block
\begin{equation}
  U(\boldsymbol{\theta}) =
  \bigl(R_Y(\theta_5) \otimes R_Y(\theta_6)\bigr)\,
  \mathrm{CNOT}_{21}\,
  \bigl(R_Z(\theta_3) \otimes R_Z(\theta_4)\bigr)\,
  \mathrm{CNOT}_{12}\,
  \bigl(R_Y(\theta_1) \otimes R_Y(\theta_2)\bigr)
\end{equation}
to disjoint neighbouring pairs, then pools by retaining every second wire. The
same $\boldsymbol{\theta} \in \mathbb{R}^{6}$ is reused across all pairs within
a stage, so the circuit carries $6L$ parameters in total --- 18 at $n=8$,
$L=3$ --- independently of the register width. The layer output is the
expectation $\langle Z \rangle$ of the single surviving wire.

\paragraph{Second ansatz family.}
A result on one circuit family is a result on one circuit family. The
trained-versus-frozen contrast is therefore repeated on a circuit that differs
from the QCNN in every respect a reviewer could name: the encoder output is
re-uploaded as $R_Y$ rotations before each of $L = 3$ layers; each layer
applies $R_Y(\theta)\,R_Z(\phi)$ to every wire, with its own parameters, and
entangles the register with a ring of CNOTs; there is no pooling; and the
output is the mean of $\langle Z \rangle$ over all eight wires. It carries
$2nL = 48$ parameters, re-uploading raises the Fourier degree of the function
it computes, and the observable is global rather than local. It is run in the
\textsc{Trained} and \textsc{Frozen} conditions on the Tokyo panel, with the
same seeds, folds, encoder, head and training loop, and compared with the
same \textsc{None} control.

\paragraph{Head.}
A two-hidden-layer perceptron consumes the concatenation
$[\langle Z \rangle, \mathbf{e}]$ and emits a logit. Training minimises binary
cross-entropy against cross-sectional extreme labels
(Section~\ref{sec:protocol:labels}).

\subsection{Batched circuit evaluation}
\label{sec:method:batched}
Evaluating the QNode once per sample, as the reference implementation does, is
the dominant cost of training. Passing the batch through the simulator as a
broadcast parameter dimension yields bit-identical expectation values --- the
maximum absolute difference over every configuration we timed is exactly zero
in double precision --- while replacing the Python-level loop. At $n=8$ and
$L=3$ the forward pass over a batch of \num{512} falls from
\SI{8.64}{\second} to \SI{0.12}{\second}, a factor of \num{71.7}; the gain
grows with batch size, from \num{33.1}$\times$ at \num{64} through
\num{64.2}$\times$ at \num{256}, because the fixed per-call overhead is
amortised over more samples (Table~\ref{tab:batching}). Adjoint differentiation
on a state-vector backend is slower at this width because the register is too
small to amortise its own overhead.

\begin{table}[t]
  \centering
  \caption{Batched against per-sample circuit evaluation at $n=8$, $L=3$,
  \num{18} parameters. Median of three repeats after warm-up; the machine was
  under scheduler control. Outputs agree exactly in double precision.}
  \label{tab:batching}
  \begin{tabular}{rrrrr}
    \toprule
    Batch & Per-sample loop (s) & Batched (s) & Speed-up & $\max|\Delta|$ \\
    \midrule
    64  & 0.765 & 0.023 & 33.1$\times$ & 0 \\
    256 & 3.218 & 0.050 & 64.2$\times$ & 0 \\
    512 & 8.635 & 0.120 & 71.7$\times$ & 0 \\
    \bottomrule
  \end{tabular}
\end{table}

This is not merely an engineering note. The reference implementation compensated
for its cost by subsampling to 9\% of trading days and by capping training at
\num{10000} sequences, which is what reduced its reported out-of-sample windows
to 22 and in some runs 2 trading days. The speed-up is what makes the evaluation
protocol of Section~\ref{sec:protocol} affordable.

\subsection{Ablation conditions}
\label{sec:method:ablation}
A hybrid model whose head also observes the classical encoding cannot, on its
own, tell us whether the circuit contributed. We therefore train four conditions
through an identical code path, differing only in the block that occupies the
quantum position:

\begin{description}
  \item[\textsc{Trained}] the circuit parameters $\boldsymbol{\theta}$ are
    learned jointly with the encoder and head.
  \item[\textsc{Frozen}] the identical circuit, held at its random
    initialisation. The contrast against \textsc{Trained} separates the value of
    \emph{optimising} the circuit from the value of the fixed nonlinear feature
    it computes.
  \item[\textsc{Surrogate}] a classical map with the same parameter count and
    the same pairwise-mixing-then-pooling structure, with angles initialised on
    $(-\pi, \pi]$ so that its output distribution matches the circuit's in
    scale (mean absolute activation \num{0.59} against \num{0.69}). A
    small-normal initialisation, the obvious default, shrinks its output by
    three orders of magnitude across three layers and would handicap the
    control.
  \item[\textsc{None}] the head receives $\mathbf{e}$ alone.
\end{description}

A fifth condition, \textsc{QuantumOnly}, withholds $\mathbf{e}$ from the head so
that the score depends on the circuit output alone. It bounds how much ranking
information the measured observable carries by itself
(Section~\ref{sec:results:ablation}).

Parameter counts across the conditions differ by less than 0.4\%
(\numrange{23817}{23899}), and the purely classical convolutional baseline is
matched to the \textsc{None} condition exactly.

\paragraph{Classical baselines.}
Five models somebody would deploy, trained on the same windows, labels,
subsample and portfolio rule: a wider temporal CNN (hidden width \num{128},
three blocks), an LSTM, a Transformer encoder, a perceptron on the flattened
window, and a gradient-boosted tree (LightGBM, \num{63} leaves, learning rate
\num{0.03}, up to \num{2000} trees with early stopping on validation AUC after
\num{100} rounds) on the same flattened $20 \times 21$ window. The last is the
model family that wins cross-sectional equity competitions, including the one
this dataset comes from, so its absence would be the first thing a finance
reader noticed.
